% Notes of the Montreal seminar, 1970 (cote 129, typed part): archive cote 129, site page 66 (PDF page 67).
% typescript. Draft compared: none.
% AI-transcribed by Claude Opus 5.5 (claude-opus-5-5, adaptive thinking, effort high), Message Batches API, batch mtl129_full (msgbatch_01AGnDKHwy6ATCM7ehUa6263), 2026-11-13; not yet checked by eye.
\scanpage{129/66}
\mspage{\hand{III-11}}

\[ D^{*}: [BT(k)]^{o} \xrightarrow{\ \approx\ } ( \hand{$\mathcal{D}$}\text{-modules, libres et de type fini sur } W) \]

Telle que
\[ [D^{*}(G)]^{\sigma} \simeq D^{*}(G^{(p)}) ; \quad D^{*}(\hand{F}_{G/k}) = F_{D^{*}(G)} ; \quad D^{*}(\hand{V}_{G}) = V_{D^{*}(G)} \]

Soit $G$ un groupe de BT sur un corps parfait. On pose
\[ M(n) = D^{*}(G(n)), \]
où $D^{*}$ est le foncteur ``classique'' défini en II, 4. Comme $G(n)$ est
annulé par $p^{n}$, $M(n)$ est annulé par $p^{n}$ et c'est donc un module sur
$W/p^{n}W = W_{n}$ .

Pour tout $i$, $1 \leq i \leq n-1$, les épimorphismes $G(n) \xrightarrow{p^{i}} G(n-i)$
définissent des monomorphismes
\[ M(n)/p^{n-i}M(n) = M(n-i) \xrightarrow{p^{i}} M(n) . \]
Ceci équivaut au fait que $M(n)$ est \uline{plat sur} $W_{n}$ (démonstration analogue
à celle de 2.2).

Enfin, le système co-p-adique des $G(n)$ définit un système
p-adique des modules $M(n)$: Les injections $G(n) \rightarrow G(n+1)$ définissent
des surjections
\[ M(n+1) \rightarrow M(n) \]
qui identifient $M(n)$ à $M(n+1)/p^{n+1}M(n)$.

Par définition, on pose:
\[ D^{*}(G) = \varprojlim M(n) \]

\mspage{\hand{66}}
